% RQIR Paper I -- CQG-targeted working manuscript
% Primary journal: Classical and Quantum Gravity (Research Paper)
% Scientific scope remains frozen at Iteration 078.
% Do not treat this working file as submission-final until CQG-I-01...09 pass.

\documentclass[12pt]{article}

\usepackage[T1]{fontenc}
\usepackage[utf8]{inputenc}
\usepackage{lmodern}
\usepackage{amsmath,amssymb,amsthm,bm}
\usepackage{booktabs}
\usepackage{graphicx}
\usepackage{microtype}
\usepackage{hyperref}
\usepackage[margin=1in]{geometry}

\newtheorem{proposition}{Proposition}
\newcommand{\Tr}{\operatorname{Tr}}
\newcommand{\RQIR}{\mathrm{RQIR}}

\title{Relativity--Quantum Interface Reconstruction I:\\
Operational source hierarchy, finite calibration, and ordered response}

\author{Aleksey Buyanov\\
Independent Researcher, Moscow, Russia\\
\texttt{pppuu7@gmail.com}\\
ORCID: \href{https://orcid.org/0009-0001-2621-9305}{0009-0001-2621-9305}}

\date{CQG-targeted working manuscript v0.3 / 8 September 2026}

\begin{document}
\maketitle

\begin{abstract}
Operational probes of the gravity--quantum interface need not be informationally complete with respect to the quantum source. We formulate a source-side reconstruction hierarchy that separates mean stress--energy data, symmetrized fluctuations, ordered response information, higher cumulants and quantum-information structure without assuming a preferred microscopic theory of gravity. At second order, the greater and lesser correlators decompose into a common symmetrized noise coordinate and an antisymmetric ordering coordinate, so a finite calibration that constrains means and noise need not determine response-sensitive information. We formalize this with a calibration-nullspace proposition: if a physical finite calibration map has a response-sensitive null direction through an interior state, sufficiently small positive density-operator perturbations can remain calibration-identical while differing in the specified ordered response. A frozen finite-dimensional RQIR construction supplies the corresponding rank, positivity and response-separation certificate, while calibration perturbations obey a steering relation controlled by the smallest retained singular value. The result is an operational non-identifiability statement. It is not evidence that gravity couples to the response sector, and it does not by itself imply quantized geometry. This source-level layer provides the finite-discriminator input for later nuisance-profiled statistical and experimental-resource analyses.
\end{abstract}

\noindent\textbf{Keywords:} quantum gravity; semiclassical gravity; stochastic gravity; stress--energy correlations; response theory; operational reconstruction; calibration; identifiability

\section{Introduction}

The experimental boundary between gravitation and quantum theory is increasingly approached through laboratory observables rather than through direct access to a microscopic gravitational state. Matter-wave and spin interferometry, gravity-mediated correlation or entanglement proposals, quantum-clock protocols, optomechanics, decoherence and diffusion tests, and dynamical field-sensitive measurements probe different projections of a common interface. For gravitational physics, the interpretive problem is therefore an inverse one: which properties of the source--gravity interface are actually fixed by a finite set of calibrated observables?

This question is important because conceptually distinct models can overlap in observable space. Semiclassical gravity, stochastic-gravity extensions, measurement-feedback or hybrid descriptions, postquantum classical-gravity models and quantum-gravitational descriptions can agree on selected low-order observables while differing elsewhere. Conversely, a nonclassical-looking detector signature does not uniquely determine gravitational ontology unless the admitted comparator class and the full source--detector map are specified.

Relativity--Quantum Interface Reconstruction (RQIR) starts from this operational ambiguity. Schematically,
\begin{equation}
 P_{\rm data}(\bm{o}\mid\bm{s})\longrightarrow [\mathfrak I],
\end{equation}
where $[\mathfrak I]$ denotes an equivalence class of interface maps after the declared preparation, calibration, baseline and nuisance freedoms are taken into account. RQIR does not assume at the outset that gravity is classical, stochastic, quantum, emergent or hybrid.

Paper I isolates the source-information layer. The central question is deliberately narrower than ``is gravity quantum?'': what source information can remain distinct after a declared finite calibration? Established closed-time-path, linear-response and stochastic-gravity structures already distinguish the stress--energy expectation value, symmetrized fluctuations and causal response \cite{FeynmanVernon1963,Keldysh1965,Kubo1957,HuVerdaguer2008}. We do not claim that separation as new. The RQIR step is to place those coordinates inside a finite physical calibration quotient and ask whether the unresolved direction can remain response-sensitive.

The main structural result is a positivity-preserving calibration-nullspace proposition. If the calibration map leaves a physical null direction and the designated ordered-response functional is nonzero along that direction, then two positive source states can be exactly calibration-equivalent while remaining response-distinct. The underlying nullspace mathematics is elementary and related to informational incompleteness in quantum tomography \cite{DArianoParisSacchi2003}; the RQIR-specific content is the gravity--quantum source interpretation, physical-state admissibility, constructive realization, detector-facing response coordinate and calibration steering.

The interpretation is intentionally conservative. We do not infer that gravity couples to the matter commutator or retarded stress--energy response. A nonzero matter commutator is not evidence for an operator-valued metric, and entanglement is not treated as an assumption-free proof of quantum gravity. The result instead identifies a source-information direction that later stages must test for statistical identifiability, physical resource cost and comparator uniqueness.

\section{Relation to established gravitational approaches}

\subsection{Semiclassical and stochastic gravity}

Semiclassical gravity uses the renormalized expectation value of the stress--energy tensor as the matter source of a classical metric. Schematically,
\begin{equation}
 G_{\mu\nu}+\Lambda g_{\mu\nu}+H^{\rm EFT}_{\mu\nu}
 =8\pi G\,\langle \hat T_{\mu\nu}\rangle_{\rm ren}.
\end{equation}
Stochastic gravity extends this mean-field description by retaining stress--energy fluctuations through the noise kernel and, in closed-time-path/influence-functional formulations, the associated response and dissipation structure \cite{HuVerdaguer2008}. A fluctuation-sensitive signal therefore cannot be identified with intrinsically quantum gravitational fluctuations without a comparator analysis.

\subsection{Information-theoretic witnesses and enlarged comparator classes}

Gravity-mediated entanglement and related information-theoretic witnesses provide a prominent laboratory route \cite{BoseEtAl2017,MarlettoVedral2017,MarlettoVedral2020,MarlettoVedral2025RMP}. Other proposals test quantumness through LOCC-type bounds without requiring entanglement generation \cite{LamiPedernalesPlenio2024}. At the same time, classical--quantum coupling, postquantum, stochastic and measurement-feedback constructions enlarge the comparator space \cite{Oppenheim2023,GalleyGiacominiSelby2023,OppenheimEtAl2023,Diosi2025,MikiEtAl2025}.

Recent debate further illustrates the dependence of witness interpretation on the admitted dynamical class. Claims that classical-gravity settings can mediate entanglement in full-QFT matter descriptions and subsequent analyses of the assumptions behind such channels motivate a comparator-explicit formulation rather than a single-witness definition of gravitational quantumness \cite{AzizHowl2025,FengVedralMarletto2026}. Dynamical probes beyond a static Newtonian potential provide a complementary reason to retain response-sensitive information in the operational hierarchy \cite{ChenGiacomini2025}.

\section{Operational reconstruction of the interface}

Let $\mathcal O=\{O_A\}$ denote operational observables. The index $A$ includes preparation, smearing or coarse graining, detector transfer, experimental settings and any renormalization prescription required to define the observable. Relative to an explicit baseline $B$ in domain $\mathcal D$, define
\begin{equation}
 \Delta_A^{(B)}(\mathcal D)
 =O_A^{\rm obs}(\mathcal D)-O_A^{(B)}(\mathcal D).
\end{equation}
Residuals defined relative to distinct baselines are not merged without an explicit transformation.

Abstractly,
\begin{equation}
 \mathfrak I:(\rho_{\rm matter},\mathcal G,\mathcal C,\lambda)
 \longrightarrow P(\bm{o}\mid\bm{s}),
\end{equation}
where $\mathcal G$ denotes gravitational/interface degrees of freedom or effective variables, $\mathcal C$ denotes constraints and consistency structure, and $\lambda$ collects model and nuisance parameters.

\section{Ordered stress--energy information}

Define
\begin{equation}
 \delta\hat T_A(x)=\hat T_A(x)-\langle\hat T_A(x)\rangle,
 \qquad
 J_A(x)=\langle\hat T_A(x)\rangle.
\end{equation}
The greater and lesser correlators are
\begin{align}
 G^>_{AB}(x,y)&=\langle\delta\hat T_A(x)\delta\hat T_B(y)\rangle,\\
 G^<_{AB}(x,y)&=\langle\delta\hat T_B(y)\delta\hat T_A(x)\rangle.
\end{align}
Introduce the symmetrized/noise coordinate
\begin{equation}
 N_{AB}(x,y)=\frac12\langle\{\delta\hat T_A(x),\delta\hat T_B(y)\}\rangle
 =\frac12(G^>+G^<),
\end{equation}
and the antisymmetric ordering coordinate
\begin{equation}
 D_{AB}(x,y)=\frac{1}{2i}\langle[\delta\hat T_A(x),\delta\hat T_B(y)]\rangle.
\end{equation}
Thus
\begin{equation}
 G^>=N+iD,\qquad G^<=N-iD.
\end{equation}
With the sign convention used here, the retarded response is
\begin{equation}
 \chi^R_{AB}(x,y)=\frac{i}{\hbar}\theta(x^0-y^0)
 \langle[\hat T_A(x),\hat T_B(y)]\rangle.
\end{equation}

The ordering-preservation rule is simple but consequential: matching $J$ and $N$ does not by itself imply matching $D$ or $\chi^R$. Such an implication requires additional physical structure---for example commutativity, spacelike separation, an equilibrium fluctuation--dissipation relation or another controlled limit. Fluctuation--dissipation relations are therefore compatible with the hierarchy; they are regime-dependent physical constraints rather than a universal identity between arbitrary calibrated states \cite{Kubo1957}.

A compact parent object is the Schwinger--Keldysh/closed-time-path generating functional
\begin{equation}
 Z_T[J_+,J_-]=\Tr\!\left(U[J_+]\rho_T U[J_-]^\dagger\right),
 \qquad W_T=-i\hbar\ln Z_T,
\end{equation}
whose derivatives generate branch-ordered correlators and response structures \cite{FeynmanVernon1963,Keldysh1965}. RQIR uses this established formalism as an organizational parent, not as a claimed novelty.

\section{Finite calibration as a physical quotient}

Let $V$ be the real tangent space of Hermitian source-state perturbations after exact hard linear preparation constraints have been eliminated. Let
\begin{equation}
 A:V\rightarrow\mathbb R^m
\end{equation}
be the declared finite calibration map and
\begin{equation}
 c:V\rightarrow\mathbb R
\end{equation}
be a specified ordered-response functional. Calibration identifies perturbations that differ by $\ker A$. If $c$ vanishes on that nullspace, the response is fixed by calibration. If a physically admissible $n\in\ker A$ satisfies $c(n)\neq0$, the response remains unresolved inside the calibration-equivalence class.

\begin{proposition}[RQIR-THM-001]
Let $V$ be the real tangent space of Hermitian source-state perturbations after all exact hard linear constraints in the declared preparation domain have been eliminated. Let $A:V\to\mathbb R^m$ be a finite linear calibration map with one-dimensional nullspace $\ker A=\mathrm{span}\{n\}$, and let $c:V\to\mathbb R$ be a linear response functional. Suppose that $c(n)\neq0$, the nominal density operator $\rho_0$ lies in the interior of the physical state set on the retained finite-dimensional Hilbert space, and $n$ is represented by a Hermitian traceless perturbation compatible with the exact preparation constraints. Then there exists $\epsilon>0$ such that
\begin{equation}
 \rho_\pm=\rho_0\pm\epsilon n
\end{equation}
are physical states, have identical declared calibration data and differ in the specified response.
\end{proposition}

\begin{proof}
Because $\rho_0$ is strictly positive, choose
\begin{equation}
 0<\epsilon<\frac{\lambda_{\min}(\rho_0)}{\lVert n\rVert_{\rm op}}.
\end{equation}
Weyl's eigenvalue bound guarantees positivity of $\rho_0\pm\epsilon n$. Trace and the exact linear preparation constraints are preserved because $n$ lies in the reduced tangent space. Since $An=0$, all declared calibration rows agree for the pair. Since $c(n)\neq0$,
\begin{equation}
 c(\rho_+-\rho_-)=2\epsilon c(n)\neq0.
\end{equation}
\end{proof}

The generic existence of null directions for incomplete measurements is not new. The application-specific statement here is that a positive source-state pair can be embedded inside the declared calibration quotient while retaining a designated ordered-response difference relevant to the gravity--quantum interface.

\section{Constructive realization and calibration steering}

The frozen Paper-I construction reports a calibration rank of $24/25$, a one-dimensional hidden direction, positive state pairs, calibration-equality residuals at approximately the $10^{-15}$ numerical scale, equal selected mean/noise coordinates, nonzero opposite ordered response and a detector-aware response coordinate. These numbers remain scientifically frozen under the Paper-I closure, but they are not yet publication-final: every one must be bound to the canonical Toy009/Toy010 generator/result authority and independently reproduced before CQG submission. No missing path will be reconstructed from manuscript prose or memory.

Calibration is also a design variable. For a smooth calibration map with normalized null vector $n$, the frozen perturbative steering relation is
\begin{equation}
 n'=-A^+A'n,
\end{equation}
with conditioning bound
\begin{equation}
 \lVert n'\rVert\leq\frac{\lVert A'\rVert}{s_{\min}(A)}
\end{equation}
in the stated setting. Changing the calibration geometry can therefore rotate the unresolved direction and alter $c(n)$ without changing the underlying source Hamiltonian. This turns calibration geometry from passive bookkeeping into an operational co-design variable.

\section{Operational sensitivity hierarchy}

The source-side hierarchy used by RQIR is summarized in Table~\ref{tab:hierarchy}.

\begin{table}[ht]
\centering
\caption{Operational source-information hierarchy. The levels classify transmitted information, not microscopic gravitational ontology.}
\label{tab:hierarchy}
\begin{tabular}{@{}lll@{}}
\toprule
Level & Independent information & Interpretation \\
\midrule
L1 & $J=\langle T\rangle$ & mean-sensitive \\
L2 & $N$ & symmetrized-fluctuation-sensitive \\
L3 & $D$ or $\chi^R$ & order/response-sensitive \\
L4 & higher CTP cumulants/responses & higher-order-sensitive \\
L5 & entanglement/channel/process data & quantum-information-sensitive \\
\bottomrule
\end{tabular}
\end{table}

Agreement at one level does not imply agreement at a higher level unless a specified physical identity supplies the implication. In particular, agreement with $\langle T_{\mu\nu}\rangle$ tests a mean-source coordinate only; a fluctuation excess can overlap stochastic-gravity, matter, detector or environmental channels; and a matter commutator is not by itself evidence for an operator-valued metric.

Low-energy gravitational EFT supplies an important controlled quantum comparator without selecting a UV completion \cite{Donoghue1994}. It belongs to the later RQIR comparator funnel rather than serving as automatic evidence that a new theory is required.

\section{From exact separation to experimental inference}

Paper I establishes an exact source-side question,
\begin{equation}
 An=0,\qquad c(n)\neq0.
\end{equation}
This does not establish that a finite noisy experiment can infer the corresponding amplitude after source, apparatus and detector uncertainties are introduced. Paper II therefore studies the nuisance-profiled information
\begin{equation}
 F_{\beta\mid\theta}=F_{\beta\beta}
 -F_{\beta\theta}F_{\theta\theta}^{-1}F_{\theta\beta},
\end{equation}
or equivalently in whitened geometry,
\begin{equation}
 F_{\beta\mid\theta}=\lVert(I-P_J)\tilde s\rVert^2.
\end{equation}
Exact source distinguishability is thus necessary for the intended discriminator but not sufficient for detector-level inferability.

A related structural warning arises when a compressed on-shell observable does not determine an off-shell/source-completed response. Schematically, additions proportional to an inverse propagator can vanish on shell while changing a source-completed derivative. Detailed source/Ward/comparator calculations belong to Paper IV; here the point is only that an observable quotient can remove information relevant to response reconstruction.

\section{Discussion}

The wider gravitational-physics significance of the construction is not that a new matter correlator has been discovered. Mean sources, noise kernels, commutators, retarded response and closed-time-path orderings are established structures. The RQIR contribution is to make their informational role explicit inside a finite physical calibration quotient and to require comparator and ontology claims to be made only after the complete operational map is specified.

RQIR-THM-001 records the minimal geometry needed for a calibration-equivalent positive pair to remain ordered-response-distinct. The finite construction shows how rank, positivity and detector-facing response can coexist, while the steering relation exposes the unresolved direction as a calibration-design variable. The result therefore supplies a concrete finite-discriminator layer that is upstream of statistical significance or apparatus feasibility.

The epistemic staging matters. Paper I identifies source-information directions. Paper II asks whether they survive nuisance profiling. Paper III converts statistically surviving directions into resource and experimental-architecture requirements. Paper IV compares existing gravity and quantum-gravity realizations through a common operational funnel. Paper V is conditional and is justified only if that comparator analysis produces a robust residual that requires a genuinely new parent model.

\section{Conclusion}

We formulated the first layer of RQIR as a calibration-aware source-information problem for the gravity--quantum interface. At second order, symmetrized fluctuations and ordering/response information are distinct coordinates unless additional physics identifies them. A finite calibration defines an equivalence class on physical source states. When that map leaves a response-sensitive physical null direction through an interior state, positive calibration-equivalent but response-distinct states exist.

This is an operational non-identifiability result, not evidence that gravity transmits the response coordinate and not evidence for quantized geometry. Its role is to specify what a finite gravitational interface experiment would have to distinguish before stronger statistical, resource or microscopic claims become justified.

\section*{Data and code availability}

The RQIR project repository contains the manuscript, analytic development and computational certificates used during the research programme. Submission-final data/code provenance is not yet claimed for the Toy009/Toy010 numerical statements: the primary canonical generator/result paths and an independent clean reproduction must be bound in the Paper-I reproduction manifest before submission. The current public project repository is \url{https://github.com/pppuu7-cmd/Relativity-Quantum-Interface-Reconstruction}.

\section*{Submission-status note}

This file is a CQG-targeted working manuscript, not a submitted or accepted article. Funding and conflict-of-interest declarations remain unspecified in the canonical author metadata and must not be inferred. Final figures, tables, bibliography normalization, Toy009/Toy010 provenance and the live CQG submission requirements remain open gates under `docs/PAPER_I_CQG_SUBMISSION_GATE.md`.

\bibliographystyle{unsrt}
\bibliography{../../docs/PAPER_I_REFERENCES_WORKING}

\end{document}
